%! Random Sampling and Data Collection — Unit 3, Lesson 3.3 — Student Packet Cover Sheet
\documentclass[10pt]{article}
\usepackage{apstats-article}
\usepackage{apstats-boxes}
\usepackage{ltablex}
\keepXColumns

\begin{document}

% Full-bleed navy banner
\begin{tikzpicture}[remember picture, overlay]
  \fill[navy] (current page.north west)
    rectangle ([yshift=-0.9in]current page.north east);
\end{tikzpicture}
\vspace{-0.8in}

\begin{center}
  {\color{white}\textbf{\LARGE Statistics}}\\[4pt]
  {\color{skymid}\large\textbf{Unit 3: Collecting Data}}\\[6pt]
  {\color{white}\textbf{\large Lesson 3.3 \quad Random Sampling and Data Collection}}
\end{center}

\vspace{0.22in}
\namedateperiod
\vspace{0.18in}

\begin{learningtargetbox}
\small
\begin{enumerate}[leftmargin=1.5em, itemsep=5pt]
  \item \ldots{} identify and describe simple random, stratified, cluster, and systematic sampling methods.
  \item \ldots{} explain why a particular sampling method is or is not appropriate for a given situation.
  \item \ldots{} distinguish stratified from cluster sampling.
\end{enumerate}
\end{learningtargetbox}

\vspace{0.14in}

\begin{tocbox}
% A CONTENTS LIST ONLY: no Score column and no Total row. Whether a component counts is said in
% its description.
\small
\renewcommand{\arraystretch}{1.5}
\begin{tabularx}{\linewidth}{@{} c l X @{}}
\rowcolor{sky}
\textbf{\#} & \textbf{Component} & \textbf{Description} \\
\midrule
1 & Warm-Up & Spiral review + sampling method preview \\
2 & Guided Notes & SRS, stratified, cluster, systematic, census \\
3 & Group Activity & Classifying and evaluating sampling scenarios \\
4 & Exit Ticket & Cluster sample of county voters \\
5 & Homework & Four sampling scenarios; classify and justify \\
\end{tabularx}
\end{tocbox}

\end{document}
